% ─────────────────────────────────────────────────────────────────────────────
%
% Documentación del mapa de ocupación 2D acumulado.
% Corresponde a la quinta iteración del módulo de percepción local.
% ─────────────────────────────────────────────────────────────────────────────

\section{Mapa de ocupación 2D acumulado}
\label{sec:lm-occupancy}

Las cuatro iteraciones anteriores del módulo de percepción local establecieron
un pipeline capaz de proyectar la imagen de profundidad a una nube de puntos
3D, alinear esa nube a la gravedad mediante odometría visual-inercial y
clasificar los puntos en suelo, obstáculo y techo. Sin embargo, este
procesamiento opera exclusivamente sobre el fotograma actual: los obstáculos
que salen del campo visual de la cámara, ya sea porque el usuario giró la
cabeza o porque el obstáculo quedó a sus espaldas, desaparecen del mapa
instantáneamente. Para un sistema de asistencia a personas con discapacidad
visual, esta carencia implica que no puede advertirse sobre obstáculos
conocidos que se aproximan desde fuera del campo de visión.

La quinta iteración introduce la clase \texttt{OccupancyMapper}, que mantiene
una rejilla 2D de ocupación acumulada en el tiempo. La rejilla es una
\textbf{ventana deslizante en el marco \texttt{odom}} (absoluto, fijo en el
mundo): los obstáculos se insertan y acumulan en coordenadas del mundo,
independientemente de la orientación de la cámara. Esto resuelve un problema
fundamental de las representaciones relativas a la cámara: cuando la cámara
gira en yaw, los ejes locales del frame relativo rotan, pero el contenido
histórico de la rejilla no puede rotarse con ellos, lo que produce un mapa
inconsistente tras cualquier giro. Al acumular en \texttt{odom}, los
obstáculos quedan fijos en el mundo y el mapa es correcto
independientemente de la orientación de la cámara.

Cuando la cámara se aleja del centro de la ventana, el contenido acumulado
se desplaza la cantidad de celdas correspondiente al movimiento en
\texttt{odom} y el origen de la rejilla se actualiza, de modo que
\texttt{worldToCell()} siga siendo correcto. Los obstáculos que caen fuera
del área cubierta se descartan.

Cada celda almacena una estimación probabilística de si esa región del suelo
está ocupada por un obstáculo. La estimación se actualiza en cada fotograma
mediante el filtro bayesiano en forma logarítmica propuesto por Elfes
\cite{Elfes1989}, y las regiones del espacio libre son marcadas mediante la
trazada de rayos con el algoritmo de Bresenham \cite{Bresenham1965}. El
desplazamiento de la ventana se calcula a partir de la odometría de
\texttt{nav\_odometry} (véase la sección~\ref{sec:lm-vio}).

El resultado se publica como \texttt{nav\_msgs/OccupancyGrid} en el topic
\texttt{/local\_mapper/occupancy\_grid} con el frame \texttt{odom},
compatible con la convención estándar de ROS 2 y visualizable directamente
en RViz.

\subsection{Marco teórico}
\label{subsec:lm-v5-teoria}

\subsubsection{Modelo de rejilla de ocupación}
\label{subsubsec:lm-v5-elfes}

El modelo de rejilla de ocupación (\emph{occupancy grid}) fue introducido por
Elfes \cite{Elfes1989}. El espacio 2D se divide en celdas de tamaño fijo
$\Delta$; a cada celda $m_i$ se le asigna una variable aleatoria binaria que
vale $1$ si la celda está ocupada y $0$ si está libre. El estado del mapa
completo es el conjunto $\mathcal{M} = \{m_i\}$.

El sensor proporciona una observación $z_t$ en el instante $t$. La estimación
bayesiana recursiva mantiene la distribución posterior:
%
\begin{equation}
  \label{eq:bayes-occ}
  P(m_i \mid z_{1:t}) =
  \frac{P(z_t \mid m_i)\,P(m_i \mid z_{1:t-1})}{P(z_t \mid z_{1:t-1})}
\end{equation}
%
Esta expresión es independiente para cada celda, lo que hace al modelo
tractable: el mapa de tamaño $N$ admite $N$ actualizaciones independientes por
fotograma \cite{Thrun2005}.

\subsubsection{Representación en log-odds}
\label{subsubsec:lm-v5-logodds}

La formulación directa de la ecuación~\ref{eq:bayes-occ} es numéricamente
problemática: las probabilidades tienden a cero o a uno muy rápidamente y
la representación en punto flotante pierde precisión. La solución estándar es
expresar la estimación en unidades de log-odds \cite{Thrun2005}:
%
\begin{equation}
  \label{eq:logodds-def}
  l(m_i) = \log \frac{P(m_i = 1)}{P(m_i = 0)} = \log \frac{P(m_i = 1)}{1 - P(m_i = 1)}
\end{equation}
%
La conversión inversa es:
%
\begin{equation}
  \label{eq:logodds-inv}
  P(m_i = 1) = \frac{1}{1 + e^{-l(m_i)}}
\end{equation}
%
El prior uniforme $P(m_i) = 0{,}5$ corresponde a $l_0 = 0$, que es el valor
inicial de toda celda sin observar.

La actualización bayesiana en la representación log-odds se convierte en una
suma, lo que garantiza estabilidad numérica y eficiencia computacional. Sea
$l_{\text{occ}}$ el log-odds del modelo de sensor al observar una celda
ocupada, y $l_{\text{free}}$ el log-odds al observar una celda libre; bajo la
hipótesis de independencia entre observaciones y prior uniforme eliminado,
la actualización queda:
%
\begin{equation}
  \label{eq:logodds-update}
  l_t(m_i) =
  \begin{cases}
    l_{t-1}(m_i) + l_{\text{occ}} & \text{si la celda es observada ocupada}\\
    l_{t-1}(m_i) + l_{\text{free}} & \text{si la celda es observada libre}
  \end{cases}
\end{equation}
%
Esta derivación, presentada en la sección~9.2 de Thrun et~al.~\cite{Thrun2005},
surge de aplicar la regla de Bayes en log-odds y cancelar el prior
$l_0 = 0$, lo que elimina la resta del término de prior en la expresión
genérica. Cada celda se actualiza en tiempo $O(1)$.

Los parámetros del modelo de sensor se eligen de forma que:
$l_{\text{occ}} > |l_{\text{free}}|$, es decir, una observación ocupada tiene
más peso que una libre. Esto garantiza que un obstáculo detectado una sola vez
requiera varias observaciones de espacio libre para ser desactivado, lo que es
deseable para la seguridad de navegación \cite{Elfes1989}. Los valores por
defecto empleados, $l_{\text{occ}} = \log(0{,}9/0{,}1) \approx 2{,}197$ y
$l_{\text{free}} = \log(0{,}3/0{,}7) \approx {-0{,}847}$, implican que son
necesarias aproximadamente $2{,}6$ lecturas de espacio libre para cancelar un
único impacto de obstáculo.

Los log-odds se acotan en el intervalo $[l_{\min},\, l_{\max}]$ para evitar
saturación numérica en celdas observadas repetidamente. Con los valores por
defecto $l_{\min} = -5$ y $l_{\max} = +5$, las probabilidades extremas son
$P \approx 0{,}67\,$\% (libre confirmado) y $P \approx 99{,}3\,$\% (ocupado
confirmado).

\subsubsection{Proyección 2D sobre el plano del suelo}
\label{subsubsec:lm-v5-proyeccion}

La rejilla representa el plano horizontal del suelo en el marco \texttt{odom}.
Cada punto 3D de obstáculo $\mathbf{p}^{\text{ga}} = (p_x, p_y, p_z)$, que
el \texttt{GroundEstimator} entrega en \texttt{gravity\_aligned\_frame}, debe
convertirse a coordenadas del marco \texttt{odom} antes de insertarse en la
rejilla. Esto se hace rotando por $\mathbf{q}_{\text{yaw}}$ (la orientación
de la cámara reducida a solo yaw, en el convenio Y-abajo descrito en la
sección~\ref{sec:lm-vio}) y sumando la posición actual del sensor:
%
\begin{equation}
  \label{eq:ga-to-odom}
  \mathbf{p}^{\text{odom}} = \mathbf{q}_{\text{yaw}} \cdot \mathbf{p}^{\text{ga}}
  + (x_{\text{odom}},\; 0,\; z_{\text{odom}})
\end{equation}
%
La proyección al plano del suelo ignora la coordenada vertical $p_y^{\text{odom}}$,
obteniendo la huella (\emph{footprint}) del obstáculo en el plano XZ de
\texttt{odom}:
%
\begin{equation}
  \label{eq:ga-to-grid}
  (x_{\text{grid}},\, z_{\text{grid}}) = (p_x^{\text{odom}},\, p_z^{\text{odom}})
\end{equation}
%
El sensor se pasa también en coordenadas \texttt{odom}: $(x_{\text{odom}},\,
z_{\text{odom}})$.

\subsubsection{Ventana deslizante: \texttt{shift}}
\label{subsubsec:lm-v5-shift}

La rejilla tiene un tamaño fijo de $N \times N$ celdas. Cuando la cámara se
desplaza y se acerca a un borde de la ventana, el contenido completo de la
rejilla se desplaza para mantenerla centrada en la posición actual del sensor.
A diferencia del enfoque relativo al frame de la cámara, donde el shift debía
expresarse en coordenadas del frame local, con la rejilla en el marco
\texttt{odom} el desplazamiento en celdas se calcula directamente desde la
diferencia de posición en \texttt{odom}, sin ninguna rotación:
%
\begin{equation}
  \label{eq:odom-delta}
  \mathbf{d}_{\text{odom}} = (x_t - x_{t-1},\; 0,\; z_t - z_{t-1})
\end{equation}
%
El desplazamiento en celdas es:
%
\begin{equation}
  \label{eq:shift-cells}
  \Delta c_i = \left\lfloor \frac{-d_{\text{odom},x}}{\Delta} + \tfrac{1}{2} \right\rfloor,
  \qquad
  \Delta c_j = \left\lfloor \frac{-d_{\text{odom},z}}{\Delta} + \tfrac{1}{2} \right\rfloor
\end{equation}
%
donde el signo negativo refleja que si la cámara avanza $+X$ en
\texttt{odom}, el contenido de la rejilla debe correrse $-X$ para que los
obstáculos queden en las celdas correctas. Las celdas que salen por un borde
se descartan; las que entran por el borde opuesto se inicializan a
$l_0 = 0$ (prior uniforme).

Adicionalmente, \texttt{shift()} actualiza las coordenadas del origen de la
rejilla en \texttt{odom}:
%
\begin{equation}
  \label{eq:origin-update}
  x_{\text{orig}} \mathrel{-}= \Delta c_i \cdot \Delta,
  \qquad
  z_{\text{orig}} \mathrel{-}= \Delta c_j \cdot \Delta
\end{equation}
%
Este paso es imprescindible: sin él, \texttt{worldToCell()} usaría un origen
desactualizado al convertir coordenadas \texttt{odom} a índices de celda en
el fotograma siguiente.

\subsubsection{Discretización en rejilla}
\label{subsubsec:lm-v5-rejilla}

La rejilla es un arreglo cuadrado de $N \times N$ celdas de lado $\Delta\,\text{m}$.
Su origen $(x_{\text{orig}}, z_{\text{orig}})$ es la posición en \texttt{odom}
del vértice $(c_i, c_j) = (0, 0)$; en el arranque del nodo se centra en la
posición inicial de la cámara:
%
\begin{equation}
  \label{eq:grid-origin}
  x_{\text{orig}}^{(0)} = z_{\text{orig}}^{(0)} = -\frac{N \cdot \Delta}{2}
\end{equation}
%
Con cada llamada a \texttt{shift()} el origen se desplaza según la
ecuación~\ref{eq:origin-update}, por lo que $x_{\text{orig}}$ y
$z_{\text{orig}}$ siempre apuntan al vértice actual de la ventana en
\texttt{odom}. La conversión de coordenadas \texttt{odom} $(x, z)$ a índice
de celda $(c_i, c_j)$ es:
%
\begin{equation}
  \label{eq:world-to-cell}
  c_i = \left\lfloor \frac{x - x_{\text{orig}}}{\Delta} \right\rfloor,
  \qquad
  c_j = \left\lfloor \frac{z - z_{\text{orig}}}{\Delta} \right\rfloor
\end{equation}
%
Un punto se considera fuera del mapa si $c_i \notin [0, N)$ o
$c_j \notin [0, N)$; en ese caso se descarta sin actualizar la rejilla.
La rejilla interna se indexa con aritmética \emph{row-major} \cite{stroustrup2013cpp}:
el índice lineal de la celda $(c_i, c_j)$ es $c_i \cdot N + c_j$, donde
$c_i$ recorre el eje $X$ y $c_j$ el eje $Z$.

Con los parámetros por defecto ($\Delta = 0{,}10\,\text{m}$, $N = 100$), la
ventana cubre un área de $10\,\text{m} \times 10\,\text{m}$ centrada en la
posición actual de la cámara. La memoria se conserva solamente dentro de un
radio de 5~m; las celdas de las esquinas del soporte cuadrado quedan en estado
desconocido por encontrarse fuera de esa región circular.

\subsubsection{Ray casting por el algoritmo de Bresenham}
\label{subsubsec:lm-v5-bresenham}

Para marcar el espacio libre entre el sensor y cada obstáculo detectado se
utiliza el algoritmo de Bresenham \cite{Bresenham1965}. Dado el segmento entre
la celda del sensor $(c_{i0}, c_{j0})$ y la celda del obstáculo
$(c_{i1}, c_{j1})$, el algoritmo genera la secuencia de celdas de la rejilla
que el segmento de recta atraviesa, usando exclusivamente aritmética entera.

El algoritmo evita errores de acumulación propios de la evaluación de la recta
parametrizada, y garantiza que cada celda intermedia sea visitada exactamente
una vez. La variable de error $e$ mantiene la diferencia acumulada entre la
recta ideal y la discretización; cuando desborda se efectúa un paso en el eje
secundario. En la implementación utilizada, con $\Delta x = |c_{i1} - c_{i0}|$
y $\Delta z = |c_{j1} - c_{j0}|$:
%
\begin{equation}
  \label{eq:bresenham-recurrence}
  e_0 = \Delta x - \Delta z, \qquad
  e \leftarrow e + 2 \Delta x \;\text{ (paso en $z$)}, \quad
  e \leftarrow e - 2 \Delta z \;\text{ (paso en $x$)}
\end{equation}
%
La celda del obstáculo en el extremo del rayo \emph{no} recibe la
actualización de espacio libre: se actualiza como celda ocupada mediante la
ecuación~\ref{eq:logodds-update} con incremento $l_{\text{occ}}$.

En condiciones típicas de uso (nube de obstáculos de $\approx\!1\,000$~puntos
a distancias de $0{,}5$--$4\,\text{m}$, resolución $\Delta = 0{,}10\,\text{m}$),
cada rayo recorre entre 5 y 40 celdas, con lo que la actualización de la
rejilla completa por fotograma tiene una complejidad práctica
$O(N_{\text{obs}} \cdot r)$, donde $N_{\text{obs}}$ es el número de puntos
de obstáculo y $r$ la longitud media de los rayos en celdas.

\subsection{Implementación}
\label{subsec:lm-v5-impl}

\subsubsection{Clase \texttt{OccupancyMapper}}
\label{subsubsec:lm-v5-clase}

La clase \texttt{local\_mapper::OccupancyMapper} encapsula la rejilla y
expone la siguiente interfaz pública:
%
\begin{itemize}
  \item \textbf{Constructor}: \texttt{OccupancyMapper(const Config\&)}.
    Inicializa la rejilla de $N \times N$ flotantes a $0{,}0$ (prior uniforme).
    Calcula las coordenadas del origen mediante la
    ecuación~\ref{eq:grid-origin}. Los parámetros clave de \texttt{Config}
    son: \texttt{cell\_size\_m}, \texttt{grid\_size}, \texttt{l\_occ},
    \texttt{l\_free}, \texttt{l\_min}, \texttt{l\_max}, \texttt{max\_range\_m}
    y \texttt{forget\_radius\_m}.

  \item \textbf{Desplazamiento}: \texttt{shift(shift\_ci, shift\_cj)}.
    Corre todo el contenido de la rejilla $(\Delta c_i, \Delta c_j)$ celdas
    según las ecuaciones~\ref{eq:shift-cells}--\ref{eq:origin-update}.
    Actualiza $x_{\text{orig}}$ y $z_{\text{orig}}$ para que
    \texttt{worldToCell()} siga siendo correcto. Las celdas nuevas se
    inicializan a $0$ (prior uniforme). Se llama \emph{antes} de
    \texttt{update()} en cada fotograma.

  \item \textbf{Actualización}: \texttt{update(obstacle\_pts, sensor\_x, sensor\_z)}.
    Recibe los puntos de obstáculo en coordenadas XZ del marco \texttt{odom}
    (según la sección~\ref{subsubsec:lm-v5-proyeccion}) y la posición del
    sensor también en \texttt{odom}. Ejecuta el pipeline de actualización
    descrito a continuación.

  \item \textbf{Resultado}: \texttt{logOdds()} retorna el vector de log-odds,
    indexado \emph{row-major} con $c_i \leftrightarrow X$.

  \item \textbf{Conversión de salida}: el método estático
    \texttt{logOddsToNavMsg(l)} convierte un log-odds a un valor entero en el
    rango $[0, 100]$ compatible con \texttt{nav\_msgs/OccupancyGrid}, o a
    $-1$ para celdas sin ninguna observación (log-odds exactamente igual a
    cero).
\end{itemize}

\subsubsection{Pipeline de actualización por fotograma}
\label{subsubsec:lm-v5-pipeline}

El método \texttt{update()} ejecuta los pasos siguientes para cada fotograma:
%
\begin{enumerate}
  \item \textbf{Conversión de la posición del sensor a celda}. Si el sensor
    (posición actual de la cámara en \texttt{odom}) cae fuera de la ventana
    de la rejilla, se omite el fotograma. En condiciones normales esto no
    ocurre porque \texttt{shift()} mantiene la ventana centrada.

  \item \textbf{Borrado local} (pase $O(N^2)$).
    Antes de insertar las nuevas observaciones, se recorre toda la rejilla.
    Las celdas ya observadas (log-odds $\neq 0$) cuya distancia al sensor
    supera \texttt{forget\_radius\_m} se reinician al prior uniforme
    (log-odds $= 0$). Esto garantiza que la información del mapa sea
    estrictamente local dentro del radio especificado. Las celdas en prior
    se omiten para reducir el trabajo en la mayoría del grid no observado.

  \item \textbf{Filtro de rango} para cada punto de obstáculo. Se calcula la
    distancia horizontal al sensor; si supera \texttt{max\_range\_m} el punto
    se descarta. Esto evita que mediciones muy lejanas con mayor ruido
    contaminen el mapa \cite{IntelRealSenseDepthTuning}.

  \item \textbf{Conversión del obstáculo a celda} mediante la
    ecuación~\ref{eq:world-to-cell}. Si la celda está fuera de los límites,
    el punto se descarta.

  \item \textbf{Actualización de celda ocupada}. Se incrementa el log-odds de
    la celda del obstáculo en $l_{\text{occ}}$, acotando en $[l_{\min},\, l_{\max}]$.

  \item \textbf{Ray casting de espacio libre} (si \texttt{enable\_raycasting = true}).
    Se traza un rayo con el algoritmo de Bresenham desde la celda del sensor
    hasta la celda del obstáculo (extremo excluido) y se aplica
    $l_{\text{free}}$ a cada celda intermedia. Para evitar que múltiples
    rayos pasen por la misma celda en un único fotograma y acumulen
    $k \cdot l_{\text{free}}$, lo que borraría en un instante el historial
    acumulado, cada celda libre solo puede actualizarse una vez por llamada
    a \texttt{update()}, mediante una máscara booleana de tamaño $N^2$.
\end{enumerate}

El método trabaja íntegramente sobre el vector interno de log-odds, sin
ninguna copia adicional de datos.

\subsubsection{Integración del \texttt{OccupancyMapper} en el nodo}
\label{subsubsec:lm-v5-transform-node}

El nodo \texttt{depth\_projection\_node} ejecuta el siguiente flujo en cada
fotograma, dentro de \texttt{onDepth()}:

\begin{enumerate}
  \item \texttt{onOdom()} (llamado cuando llega un mensaje de odometría)
    almacena la posición actual de la cámara $(x_t, z_t)$ en el marco
    \texttt{odom} y el cuaternión de yaw $\mathbf{q}_{\text{yaw}}$.

  \item \texttt{onDepth()} calcula el desplazamiento de la cámara desde el
    fotograma anterior con la ecuación~\ref{eq:odom-delta} y obtiene el shift
    en celdas con la ecuación~\ref{eq:shift-cells}. Llama a
    \texttt{occupancy\_mapper\_.shift($\Delta c_i$, $\Delta c_j$)}, que
    también actualiza el origen según la ecuación~\ref{eq:origin-update}.
    Si todavía no hay posición previa registrada (primer fotograma), el
    \emph{shift} se omite.

  \item Los puntos de obstáculo obtenidos del \texttt{GroundEstimator}
    (en \texttt{gravity\_aligned\_frame}) se transforman a \texttt{odom}
    mediante la ecuación~\ref{eq:ga-to-odom}, se proyectan al plano XZ
    (ecuación~\ref{eq:ga-to-grid}) y se pasan a
    \texttt{occupancy\_mapper\_.update()} junto con la posición del sensor
    $(x_{\text{odom}}, z_{\text{odom}})$.

  \item Tras la actualización, si al menos un suscriptor escucha el topic
    \texttt{/local\_mapper/occupancy\_grid}, se construye y publica el mensaje
    \texttt{nav\_msgs/OccupancyGrid}.
\end{enumerate}

\subsubsection{Publicación del \texttt{OccupancyGrid}}
\label{subsubsec:lm-v5-occgrid}

El mensaje \texttt{nav\_msgs/OccupancyGrid} utiliza la convención de indexado
$\text{data}[\text{col} + \text{row} \cdot \text{width}]$, donde
\texttt{col} crece en la dirección $X$ y \texttt{row} crece en la dirección
$Y$ del frame del mapa. Dado que la rejilla interna indexa en \emph{row-major}
con $c_i \leftrightarrow X$ y $c_j \leftrightarrow Z$.  Dado que la
rotación de $+90^{\circ}$ en el origen alinea la dirección $X$ local con
$c_i$ y la dirección $Y$ local (filas) con $c_j$, la correspondencia
con el campo \texttt{data} es directa; véase la
ecuación~\ref{eq:grid-data-index}.
%
El campo \texttt{info.origin} del mensaje recibe las coordenadas actuales del
vértice $(0, 0)$ de la rejilla en \texttt{odom}, con
\texttt{origin.position.x} $= x_{\text{orig}}$,
\texttt{origin.position.y} $= h_{\text{piso}}$ (altura del piso estimada por
el \texttt{GroundEstimator}) y
\texttt{origin.position.z} $= z_{\text{orig}}$.

El campo \texttt{info.origin.orientation} recibe un cuaternión de $+90^{\circ}$
alrededor del eje $X$, con componentes $(w, x, y, z) = (\tfrac{\sqrt{2}}{2},
\tfrac{\sqrt{2}}{2}, 0, 0)$. Esta rotación hace que el eje $Y$ local del
grid, en el que \texttt{nav\_msgs/OccupancyGrid} extiende las filas, quede
alineado con el eje $Z$ de \texttt{odom} (dirección de avance). Sin esta
rotación, el grid se dibujaría en el plano $XY$ vertical de \texttt{odom}
en lugar del plano $XZ$ horizontal que corresponde al suelo.

Con la orientación corregida, la celda $(c_i, c_j)$ se ubica en \texttt{odom}
en:
%
\begin{equation}
  \label{eq:cell-world-pos}
  \mathbf{p}_{\text{odom}}(c_i, c_j)
  = \bigl(x_{\text{orig}} + c_i \Delta,\; h_{\text{piso}},\; z_{\text{orig}} + c_j \Delta\bigr)
\end{equation}
%
La relación entre el vector interno de log-odds y el campo \texttt{data} del
mensaje es directa, sin transposición:
%
\begin{equation}
  \label{eq:grid-data-index}
  \texttt{data}[c_i + c_j \cdot N]
  = \texttt{logOddsToNavMsg}\!\left(\text{grid}[c_i \cdot N + c_j]\right)
\end{equation}
%
donde $c_i$ corresponde a la columna (\texttt{col}, dirección $X$) y $c_j$
a la fila (\texttt{row}, dirección $Z$ tras la rotación de origen).

\subsection{Integración en el nodo \texttt{depth\_projection\_node}}
\label{subsec:lm-v5-integracion}

La Tabla~\ref{tab:v5-topics} resume el topic nuevo introducido por esta
iteración.

\begin{table}[ht]
  \centering
  \caption{Topic publicado por el \texttt{OccupancyMapper} (V5).}
  \label{tab:v5-topics}
  \begin{tabular}{llp{6cm}}
    \hline
    \textbf{Dirección} & \textbf{Topic} & \textbf{Descripción} \\
    \hline
    Pub & \texttt{/local\_mapper/occupancy\_grid} &
      \texttt{nav\_msgs/OccupancyGrid}.
      Rejilla 2D deslizante en el marco \texttt{odom}
      (QoS: \texttt{transient\_local}). \\
    \hline
  \end{tabular}
\end{table}

La clase \texttt{OccupancyMapper} se compila dentro de la librería estática
\texttt{local\_mapper\_lib}, junto con \texttt{DepthProjector},
\texttt{GroundEstimator} y las demás clases del paquete. El ejecutable
\texttt{depth\_projection\_node} enlaza con \texttt{local\_mapper\_lib} sin
ninguna dependencia adicional en tiempo de compilación.

\subsection{Parámetros configurables}
\label{subsec:lm-v5-params}

La Tabla~\ref{tab:v5-params} lista los parámetros del
\texttt{OccupancyMapper} declarados en el nodo y configurables desde
\texttt{params.yaml}.

\begin{table}[ht]
  \centering
  \caption{Parámetros del \texttt{OccupancyMapper} (prefijo \texttt{occupancy\_}).}
  \label{tab:v5-params}
  \small
  \begin{tabular}{lcp{7cm}}
    \hline
    \textbf{Parámetro} & \textbf{Valor por defecto} & \textbf{Descripción} \\
    \hline
    \texttt{cell\_size\_m}        & $0{,}10$  & Lado de cada celda cuadrada. \\
    \texttt{grid\_size}           & $200$     & Celdas por lado ($N$). Cubre $20\,\text{m}\times20\,\text{m}$. \\
    \texttt{l\_occ}               & $2{,}197$ & Incremento log-odds por hit de obstáculo. \\
    \texttt{l\_free}              & $-0{,}847$ & Incremento log-odds por celda libre (rayo). \\
    \texttt{l\_min}               & $-5{,}0$  & Saturación inferior del log-odds. \\
    \texttt{l\_max}               & $+5{,}0$  & Saturación superior del log-odds. \\
    \texttt{max\_range\_m}        & $5{,}0$   & Distancia máxima de los obstáculos incluidos (m). \\
    \texttt{forget\_radius\_m}    & $5{,}0$   & Radio de localidad (m); las celdas más alejadas se reinician al prior en cada fotograma. \\
    \texttt{enable\_raycasting}   & \texttt{true} & Activa la actualización de celdas libres. \\
    \hline
  \end{tabular}
\end{table}

\subsection{Decisiones de diseño}
\label{subsec:lm-v5-design}

\paragraph{Marco \texttt{odom} en lugar de frame relativo a la cámara.}
Una alternativa natural es acumular los obstáculos en coordenadas relativas
al frame de la cámara (\texttt{gravity\_aligned\_frame}). Sin embargo, este
enfoque tiene un defecto fundamental: cuando la cámara gira en yaw, los ejes
locales del frame relativo rotan, pero el contenido histórico de la rejilla
(un arreglo de celdas rectangulares) no puede rotarse en el mismo acto. El
resultado es que los obstáculos acumulados quedan con un ángulo incorrecto en
el nuevo sistema de referencia tras cualquier giro. Además, en el enfoque
relativo el \emph{shift} en celdas debe expresarse en el frame rotado
(aplicando $\mathbf{q}_{\text{yaw}}^{-1}$ al desplazamiento en \texttt{odom}),
lo que introduce una dependencia adicional y una fuente extra de error.

Al acumular en \texttt{odom} (frame absoluto, fijo en el mundo), los
obstáculos quedan fijos en sus coordenadas del mundo independientemente de
la orientación de la cámara. El \emph{shift} se calcula directamente desde
la diferencia de posición en \texttt{odom}, sin ningún paso de rotación.
El costo es que cada punto de obstáculo debe convertirse de
\texttt{gravity\_aligned\_frame} a \texttt{odom} antes de insertarse
(ecuación~\ref{eq:ga-to-odom}), operación de costo $O(N_{\text{obs}})$
mediante una rotación de cuaternión por punto.

\paragraph{Localidad estricta mediante \texttt{forget\_radius\_m}.}
En cada fotograma las celdas cuya distancia al sensor supera
\texttt{forget\_radius\_m} se reinician al prior uniforme. Esto mantiene el
mapa estrictamente local: no persisten datos a distancias a las que el sensor
no puede detectar obstáculos de modo confiable (acorde con
\texttt{max\_range\_m}), y la memoria del grid no acumula err or de odometría
de forma indefinida. El radio de olvido se configura con el mismo valor
por defecto que el rango máximo de detección ($5\,\text{m}$).

\paragraph{Sin dependencias externas.}
La clase \texttt{OccupancyMapper} está implementada en C++17 puro, sin
dependencias de PCL, Eigen ni ninguna otra librería externa. Esto es
consistente con el resto del paquete \texttt{local\_mapper} y facilita el
testing unitario y la ejecución en la Raspberry Pi 5.

\paragraph{Orientación del origen y mapeo de índices.}
La indexación \emph{row-major} interna con $c_i \leftrightarrow X$ y
$c_j \leftrightarrow Z$ es práctica para el ray casting y la actualización
de celdas. El campo \texttt{data} de \texttt{nav\_msgs/OccupancyGrid}
indexa primero en la dirección local $X$ (columnas, eje rápido) y luego
en la dirección local $Y$ (filas, eje lento). La rotación de $+900^{\circ}$
alrededor de $X$ aplicada en \texttt{info.origin.orientation} hace que la
dirección local $Y$ del grid coincida con el eje $Z$ del marco
\texttt{odom}, de modo que la ecuación~\ref{eq:grid-data-index} no
requiere transposición: \texttt{data}[$c_i + c_j \cdot N$] recibe
directamente el valor de la celda ($c_i$, $c_j$) de la rejilla interna.
